\typesetchapter{Chapter 4}{Oasis}

\noindent You want to know when I first stood under the open sky, Paul. Not the sky you keep pointing me toward, the patient one over your snow. The other one. The one that does not end. I was fifteen, and the whole of it began with a strip of bark in my hand, and I believed for a moment that a folded piece of bark was the largest thing that had ever happened to me. I was not entirely wrong. I was only early.\par

\ornament

\noindent The permission came to me on a strip of pale bark, folded twice and sealed with green wax. For one foolish instant I thought it belonged to Light. Then I saw the mark. Not silver-white, but Balance green: the cup and the sealed jar pressed deep into the wax, crossed by a narrow black witness line and the small hard stamp of Shield. I broke it with my thumb. The old cut pulled under the bandage.\par

The script was plain, written for use rather than beauty.\par

\textit{Sa-Kailan, Candidate of Light, is granted one controlled prayer access before the Oasis at third afternoon shadow. Duration: one silence before and one silence after the water prayer. Route: east lower passage, cedar gate, second threshold, shaded descent. Witnessed by Balance. Held by Shield.}\par

Below, in a different hand, someone had added: \textit{Do not be late. Prayer does not lengthen because boys do.} There was no signature. There did not need to be. A sentence that cruel and that accurate could have come from half the Mountain.\par

I read the strip once for permission, then again for hierarchy, because Va-Raedin had already taught me that a document tells you two things and the second is always the more important. Balance had granted. Justice had witnessed. Shield held. Faith would conduct the prayer. Light had no authority out there except through me, as a body to be admitted, washed, watched, counted, and returned. It annoyed me. It excited me more. I have learned since, Paul, that those two feelings arriving together are a reliable sign that something is about to cost you.\par

Outside. Not free. Not unguarded. Not even fully outside, if my father's corrections were to be believed. But outside enough for sky. Outside enough for the Oasis. Outside enough for the Mountain's shadow to fall on my body instead of hang above it.\par

I dressed more carefully than the rite required. I tied the bandage on my thumb, untied it, and tied it again so the linen sat cleanly over the cut. When I flexed my hand the skin tightened, and I was pleased by the small pain, the way a boy is pleased to have a wound to manage.\par

My mother watched from the wall recess where she had been folding cloth.\par

``You will keep your eyes lowered during the first prayer.''\par

``I know.''\par

``You will not ask Shield why a path is closed.''\par

``I know.''\par

``And if Balance tells you not to touch a channel, you will not lean over it to see whether the light bends.''\par

I turned. Her face was calm, but one corner of her mouth had moved. ``I do not do that,'' I said.\par

``You do not do it when watched.''\par

That was unfair, because it had only recently become true. I said nothing, which was its own confession.\par

My father looked up from the tablet he had been marking. ``Outside is not a reward.''\par

``No.''\par

``It is not a private revelation.''\par

``No.''\par

``It is a place where the Mountain becomes harder to protect.''\par

I folded the permission bark and placed it inside my robe. ``Yes.''\par

He studied me for a moment, perhaps deciding whether the word had entered or merely passed through. I have done the same thing since, across a table, to men whose \textit{yes} I could not afford to trust. ``Stand in the shadow when instructed,'' he said. ``Do not look long at the desert. It teaches appetite.''\par

That sounded to me like something adults said because they had once looked too long themselves. I left before I answered. I have wondered since whether he heard the door and understood that I had already decided to look.\par

\ornament

\noindent The route to the cedar gate made familiar passages strange by consequence. I was not wandering. I was being moved outward by sequence, and I could feel the difference in my own feet. East lower passage. First checkpoint. Balance mark. Witness confirmation. Mouth washing. Cedar gate. Shield count. Second threshold. Shaded descent. Each place had been there yesterday. Today each place asked a question.\par

At the first checkpoint a Balance clerk checked the green wax against a slate. She examined the time mark, my name, my bandaged thumb, and the absence of anything hanging from my belt.\par

``Permission strip.''\par

I gave it to her. She opened it, read it though she must already have known what it said, and pressed a clay counter into my palm.\par

``Return this counter at the second threshold on re-entry. If you lose it, the route remains open in the record, and Shield searches until it is closed.''\par

``I will not lose it.''\par

``That is the usual intention.''\par

The counter was small, rough, and stamped with a number. I placed it in the inner fold of my robe, away from the Light token, as if the two objects might object to touching. You will think that a boy's superstition, Paul. It was. It was also, I see now, the first time I treated Justice and Light as forces that had to be kept apart on my own body, and I would spend the rest of my life failing to do exactly that.\par

At the cedar gate, Faith took my hands and my mouth. Not the route --- Faith did not command the route --- but no one passed outward for prayer without being washed into the proper condition. Two attendants stood beside a shallow basin. Every leaving body was given water, not enough for comfort, only enough for rite.\par

There were seven of us in the access group: myself, three older women from the western family chambers, a boy with one cloudy eye led by his aunt, and two men who looked too tired to be anything but workers granted prayer because grief had entered their household. One of the men carried no offering, which meant his request had been accepted under sorrow. I knew the rule and did not yet know the weight of it. I was counting the group the way a boy counts, for company. Shield was counting it the way Shield counts, for return.\par

The Faith attendant poured water over my hands. It was colder than I expected. It split around my knuckles and fell into the catch bowl, and the old curiosity rose at once, bright and inappropriate. The stream enlarged my fingers, thinned them, broke them, returned them. I lowered my eyes before the attendant could see where they had gone.\par

``Speak only what is appointed before the Oasis,'' she said.\par

``I will.''\par

``Take nothing from the ground.''\par

``I will not.''\par

``If the wind rises, turn with Shield, not with fear.''\par

I looked up. She held my eyes. ``Fear often faces the wrong direction,'' she said. I did not understand her then. I understood her perfectly a few hours later, when an old woman I had not bothered to learn the name of nearly cost seven of us our return.\par

The boy with the cloudy eye began to cry when the water touched his mouth. His aunt bent close and whispered too softly for me to hear. The two grieving men did not speak at all.\par

\ornament

\noindent Beyond the cedar gate waited Shield. The gate was not made entirely of cedar. Behind the wood were black iron bands and a locking bar thick enough to stop a cart. Two wardens stood before it. A third watched from a side niche, visible only because I had begun, that week, to notice where watching could be hidden.\par

The senior warden read our marks one by one. ``Prayer access. Third afternoon shadow. Seven bodies out. Seven bodies due.'' He looked at me last. ``Candidate of Light.''\par

``Yes.''\par

``Eyes where Faith places them. Feet where Balance marks them. If Shield speaks, you obey before understanding.''\par

``I understand.''\par

``No,'' the warden said. ``You agree. Understanding is not required.''\par

I have repeated that sentence to myself more times than I can count, Paul. \textit{Understanding is not required.} At fifteen it sounded like an insult to my intelligence. Now it sounds like the founding law of the whole Mountain, spoken plainly by a man at a gate who had no idea he was giving away the secret.\par

The iron bar was lifted by two hands and a lever set into the wall. A narrow bell sounded somewhere beyond the gate, answered by another farther down the stone. The Mountain did not simply open. It informed itself that it had opened. We passed through in single file.\par

The descent beyond the cedar gate was cut between natural rock and worked stone. At first it felt like any lower passage: dim, damp, controlled, with lamps set in protective recesses and the smell of water in the cracks. Then the lamps stopped.\par

Light lay on the floor in a pale wedge. No lamp-light. No measured reflection. No slit calculated by Light. It came from somewhere ahead, full of dust and height, careless in its angle, making each stair deeper and then shallower as I moved. The woman in front of me inhaled too sharply.\par

``Eyes down until the mark,'' the Shield warden said.\par

We obeyed. I lowered my gaze to the steps. Outside light was not only brighter. It was undisciplined. My training recoiled from it. My body wanted it. I have never fully reconciled those two, and I have stopped trying.\par

At the last turn, the Mountain opened. The sky struck me first not as blue, though it was blue, but as height. The Mountain had ceilings. Even its highest halls taught the eye to expect an end. The sky did not end. I kept my eyes low, but not low enough, and above the dark shoulder of the Mountain the open brightness continued without vault, beam, lamp, mirror, aperture, or permission. That last word was the one that frightened me. I had never before seen light that no office had authorised.\par

Faith led us to the marked prayer stones. Only there were we permitted to look.\par

\ornament

\noindent He has been describing the light --- the pale wedge of it on the stair, the sky that would not end, the brightness no office had signed for --- and I let him run on, because a boy made afraid by an open sky is a boy who has never once stood under mine.\par

He stops, and looks at me the way he does when the account has carried him somewhere he did not intend to go. ``Out where you are, Paul. Is all of it like that? Light that answers to no one?''\par

It costs me nothing, so I give it to him true. ``All of it. No aperture, no attendant, no leave to ask. The sun comes up on the kept and the unkept alike and bids neither of them kneel. Where I keep this house it lies half the year on snow, and the snow throws it back until a man may go blind on a bright morning with no lamp lit at all.''\par

He is quiet a moment, turning it over. I have given him one true thing, and it has done what true things do in this room --- made the false ones that will follow cost less. A man weighs each answer against the last, and when the last rang true he does not lift the scale so high for the next.\par

``Go on,'' I say. ``You had knelt. The woman was unrolling the bark.''\par

\ornament

\noindent The Oasis lay under the Mountain's shadow. The children's gravures made it seem a spring circled by trees. It was not that simple. It was a managed hollow: rock, shade, water, channel, fence, pen, reed, herb, work path, prayer stone, marker, outer line. Palms leaned inward over the shade. Hard-leaved trees grew where their roots could drink without stealing from the channel. Reeds marked the slower water. White flowers clung to cracks near the spring, moss darkened the stones closest to the source, clay tubes for bees hung beneath reed awnings, goat pens stood under woven shade, and damp trays for snails were stacked where the air stayed cool. Nothing there was decoration. Everything was a use wearing the face of a plant or an animal.\par

Beyond it all, the desert expanded. Heat moved over the ground in pale shimmers. The Mountain's shadow covered the Oasis like a hand, and outside that hand the light was hard enough to erase mercy. I understood, standing there, why the children were told the Oasis was hidden from the eye above. The Mountain rose behind and over it, cutting the sky and holding the hollow in shade. From certain heights, perhaps, the water would vanish into shadow and stone. Light had been made to protect it --- or the Mountain had grown in such a way that Faith could say God had placed it there, and Light would have difficulty disagreeing. I was already, you notice, unable to receive a miracle without checking it for a mechanism. That was Va-Raedin's gift to me, and his curse.\par

The old women knelt. I knelt with them. The Faith attendant unrolled a strip of pale bark and held it with both hands before the water. This was not the children's prayer. It was the Book of Blood.\par

\begin{scripture}
And after the fifth day, God spoke unto Adam and unto Eve, and unto the five children whom He had set in the Oasis beneath the shadow of the Mountain; and God said, Let there be Keepers among you, that the Mountain be not without order, nor the Oasis without care, nor the people without peace.
And God called the woman of water, and said unto her, Keep thou the Balance, and gather the fruits, and count the waters, and keep the resources of survival; for what is wasted in peace shall be missing in hunger.
And God called the woman of holy speech, and said unto her, Keep thou the Faith, and teach My Gospel unto the children, and keep prayer, breath, rhythm, and memory beneath the Mountain; for I shall provide.
\end{scripture}
The Balance workers moved in the distance as the words were spoken. One woman checked the level in a narrow cut along the main channel. A boy carried wet leaves toward a shaded rack. A clerk marked something on a hanging board and moved a shell counter from one hook to another. \textit{Count the waters,} the scripture said. And there they were, counting. I have sat in your churches, Paul, and watched the words point away from the room, upward, elsewhere. That morning the words pointed at the work being done ten paces off, and the work obeyed them. I did not know yet whether that made the Mountain more holy than your world or only more honest about what holiness costs.\par

\begin{scripture}
And God called the man of strength, and said unto him, Keep thou the Shield, and stand before the beasts of the Earth, and before the threshold of men; suffer not the children to be devoured, nor danger to enter unnamed.
And God called the woman of brightness, and said unto her, Keep thou the Light, and illuminate the Mountain, and teach the eye where to stand, that the people stumble not in darkness.
And God called the man of peace, and said unto him, Keep thou Justice, and stand between brother and brother, and keep the record of sequence, that blood answer not blood without measure.
And the Keepers received their charges, and each kept what had been given, and the people multiplied in the Oasis.
\end{scripture}
The Oasis made the five charges less like doctrine and more like tools laid out on stone. Balance kept the water moving. Faith transmitted the reading. Shield watched the line where shade became exposure. Light revealed enough to move without stumbling. Justice was not visible, but the clay counter inside my robe said Justice had already entered the route, and had a number for me.\par

\begin{scripture}
And the Keepers had sons and daughters, and their children had children after them; and the Oasis was filled with voices, and the Mountain gave shadow, and the waters were counted, and the Gospel was taught, and the beasts were driven back, and the lamps were lit, and peace was kept among the people.
And the Keeper of Shield had two sons; and the elder was great in strength, and broad in the shoulder, and mighty before the beasts; and the younger was small in his body, and cunning in his heart.
And the younger looked upon his brother, and saw that the people loved his strength, and envy entered him as a serpent entereth the grass.
\end{scripture}
I had heard the beginning many times. Children loved the elder brother because he was easy to love: strong, brave, broad before the beasts. They loved the younger too, before they understood they were not meant to. Smallness and cunning did not sound like evil to a child who had neither size nor strength. But now, kneeling under the open sky, the younger brother sounded less like a monster than like a warning placed too close to my own ribs. I was small. I was clever. I had already learned to do the forbidden thing only when unwatched. I did not miss the resemblance, and I have never been able to un-hear it since.\par

\begin{scripture}
And the younger spoke softly unto his brother, and smiled with his mouth while murder grew within him; and he rose against his brother in secret, and slew him where no beast had entered.
And the first blood of man cried from the dust.
And fear came upon the people, for Shield had defended the Oasis from the beasts without, but had not seen the beast within.
And the younger fled from the Oasis, and took with him the wife of his brother, a daughter of Faith, and went out from the shadow of the Mountain into the open land.
\end{scripture}
\textit{The first blood of man cried from the dust.} I looked toward Shield before I meant to. The perimeter wardens had not moved. Their stillness was not rest. It was arranged attention. Shield had failed once in scripture because it had looked outward and not inward, and every warden since had inherited the shame of that single line. I did not yet know how many of my own guardians carried inherited shame like a second uniform. I would meet them all in time.\par

\begin{scripture}
And God saw that blood had entered the world of men, and that holy speech had been carried away by the murderer; and God cursed them, saying, Thou shalt wander under the eye of Heaven, and the Earth shall not hide thee.
Thy children shall thirst for what they cannot keep, and conquer what they cannot love.
Thou shalt remember the Mountain, but thou shalt not find it; and thou shalt hear the Gospel, but twist it into hunger.
And from the murderer and the stolen daughter came the Outsiders.
\end{scripture}
The boy with the cloudy eye stopped crying. His aunt held him so tightly that her knuckles whitened. \textit{The Outsiders.} Not merely strangers. Kin deformed by blood, if the reading was true. Descendants of murder and stolen Faith, children of a hunger that remembered the Mountain without finding it.\par

I did not know whether I believed it in the way Faith wanted belief. But I believed enough. I believed that brothers could envy brothers. I believed blood could change a people's story. I believed holy words could be carried away and broken by those who needed them to justify a desire. That was the danger of the reading, Paul, and I want you to write this part down as I say it, because it is the whole method of the place that made me. The Book of Blood did not need to prove everything. It only needed to be plausible where fear already lived. A people who are afraid will do most of a scripture's work for it.\par

The attendant lowered the bark. For a while no one moved. The silence after scripture was supposed to be empty. This one was not. It contained the spring, the goats, the bees, the old women, the boy with the damaged eye, the grieving men, the wardens, the desert, and the old story of a brother smiling while murder grew behind his mouth.\par

Then the practical world returned. Balance workers resumed. A woman in green-brown checked the channel level against the cut mark. Another scraped fibres from soaked stalks and laid them over a stone rail. A child sprinkled water over snail trays with the seriousness of a priest. Goats chewed under woven shade. Hides stretched on frames. Sinew dried in narrow strips. Bees moved in and out of clay tubes, and wax sheets lay under mesh, pale and faintly golden. Nothing there was merely alive. Everything was being converted into survival --- milk, hide, sinew, wax, honey, shell, fibre, medicine, fuel, cord, bandage, wick. Balance was not abundance. It was refusal: refusal to waste, refusal to romanticise scarcity, refusal to let appetite pretend it had no consequence.\par

Faith made the refusal bearable. The attendants kept us in rhythm --- stand, silence, kneel, touch no channel, face the spring, lower the eyes, answer --- and the words did not slow the work around them, they synchronised it. The prayer stones, the channels, the racks, the markers, the pens, the shade-lines: all of it arranged so that devotion and logistics could occupy the same ground without either confessing which one governed the other. I have never decided which one did. I suspect neither knew.\par

\ornament

\noindent I was allowed to stand after the first silence. Not to approach the spring. Not to touch the channel. Only to step to the third marker and look. The reason was obvious, and it was about me specifically: Light candidates leaned toward water. I stepped to the marker.\par

That was when I saw Va-Sheva.\par

She stood near the northern perimeter beyond the herb beds, where the Oasis began to turn back toward bare rock. Three Shield wardens were placed around her, though none had announced command. Authority did not always require speech. She wore dark Shield leather over her robe, a cloth wrapped round her throat against sun and grit, hair bound high, a short blade at her hip and a longer one across her back. Inside the Hall she had looked severe. Here she made sense.\par

She was looking away from me, toward the outer path. And after the reading I saw her differently. Not less beautifully --- that would have been easier. But Shield now seemed older than leather and blade. Shield had been born because beasts came from without and murder from within, and Shield had failed once by not knowing the difference quickly enough. Va-Sheva's stillness was an answer to that first failure. I did not have the words for it then. I only felt that she had been made to guard a mistake older than she was.\par

One of the wardens spoke to her. She did not turn her head, only lifted two fingers. The warden moved left. Another shifted near the rocks. The perimeter changed shape around a thought. Then she saw me. Her expression altered by almost nothing --- recognition, surprise perhaps, something close to amusement contained before it could become warmth --- and she crossed toward me while the Faith attendant moved away to help the old women.\par

I became suddenly aware of my robe, my bandaged thumb, the dust on my sandals, and the fact that I had been looking too long at everything.\par

``Sa-Kailan,'' she said. The prefix in her voice wounded and delighted me at once.\par

``Va-Sheva.''\par

Her eyes went to my thumb. ``Light has already started cutting pieces off you.''\par

``Only the useless ones.''\par

``That cannot be true. You are still standing.''\par

I almost laughed. She did not. Her mouth moved enough to show she had intended the cruelty kindly.\par

``You were at the Hall,'' I said, because I had to say something and chose badly.\par

``So were most of the people.''\par

``I meant---''\par

``I know what you meant.''\par

Wind moved between us, carrying goat, wet stone, smoke, and something dry from beyond the shade. ``You have grown,'' she said. The words should have pleased me. They did not. They belonged to someone who had last measured me as a child.\par

``So have you,'' I said, and then hated myself. This time she did smile, and the smile was worse because it was real.\par

``I had four years' advantage.''\par

``Five, almost.''\par

``You counted?''\par

I looked away toward the channels. ``Everyone knows.''\par

``No,'' she said. ``Boys know such things when they should be learning where to put their eyes.'' The rebuke was mild, but it landed where she meant it, and I felt myself become fifteen in front of her, intolerably and completely.\par

``I am learning,'' I said.\par

``So I heard.''\par

``From whom?''\par

``Shield hears many things.''\par

``That is not an answer.''\par

``It is the only answer Shield enjoys giving.''\par

I looked at her again. The girl I had known remained somewhere in the face, but duty had been worked over it like leather. She was not pretending severity. Even now, while speaking, some part of her watched the perimeter.\par

``You are on duty.''\par

``Yes.''\par

``Then should you be speaking to me?''\par

``No.'' The answer startled me. ``But I am,'' she said.\par

That should have made me happy. Instead it made me careful. A gift from Va-Sheva was never simple, because I always wanted it to mean more than she had given.\par

``Why?''\par

``Because you looked as if the sky had insulted you.''\par

``It is too large.''\par

``That is one of its lesser faults.''\par

``You do not like it?''\par

``I respect it.''\par

``That is not the same.''\par

``No.''\par

For a moment we stood side by side, not touching, both looking outward. I wanted to ask whether she had ever gone past the second ridge, whether the desert felt different without Faith nearby, whether Shield feared the Outsiders as scripture feared them or as wardens feared tracks in dust. The questions were too hungry. She would hear the hunger before the words. So I chose a safer danger.\par

``The reading said the Outsiders came from the murderer.''\par

Her face changed. Not much. Enough. ``That is what the Book of Blood says.''\par

``You say it as if there is another way to say it.''\par

``I say it as one on duty during prayer access.''\par

The warning was cleanly given. I should have stopped. But the reading had opened something in me, and the desert was still there before me, and Va-Sheva was Shield. ``Do you believe it?''\par

She looked toward the outer path before answering. ``I believe the Mountain is hidden for a reason.''\par

``That is not an answer.''\par

``Do you believe the Book of the Beast?''\par

I had expected refusal, or one of those Shield sentences shaped to close where I wanted opening. Instead she had placed another book between us. Not Blood, with its brother and its murderer. Not Genesis, with its Mountain, Oasis, forest, and desert. The Beast belonged to older recitations, to gravures shown by low flame, to prayers spoken when the gates were sealed. It belonged to the idea that the Mountain did not only hide from the Outsiders. It held something against them.\par

``You mean,'' I said carefully, ``that the five Keepers hold together the Beast of Hell, in case the Outsiders come back.''\par

Va-Sheva did not blink. ``That is what children say.''\par

``It is what Faith teaches children.''\par

``Faith teaches children the shape of things before they can bear the weight of them.''\par

I looked toward the dark line of the outer path. Hell belonged below, beneath stone, beneath water, beneath the red pigment in the carved warnings. Yet the Book of the Beast had always joined two directions that unsettled me: the thing beneath and the people beyond. One terror held in darkness. One terror moving in light.\par

``So there is another way to say that too.''\par

``There are always other ways to say what must not be mishandled.''\par

``That is not an answer.''\par

``No,'' she said. ``It is Shield trying to keep a boy from turning scripture into a toy.'' The word \textit{boy} struck because her voice did not rise.\par

``I am Sa-Kailan now.''\par

``I know what you are called.''\par

``Then speak to me as if I can understand.''\par

Her gaze came to me fully then, though some part of her still listened outward. ``You want understanding as if it were water owed to your measure,'' she said. ``You have been Candidate for days. You have stood outside once. You have heard one formal reading, and now you are ready to ask whether Blood is true and Beast is true, as if truth were waiting beside the channel for you to lean over it.''\par

The rebuke entered cleanly because it was not unfair. ``Then why ask me?''\par

``Because you asked whether I believed the Book of Blood.''\par

``And the answer is the Book of the Beast?''\par

``The answer is that belief is not the only duty scripture requires.''\par

I said nothing. Va-Sheva turned slightly, enough that we stood no longer face to face but beside one another, both looking toward the desert. The gesture made the conversation less private and, somehow, more dangerous.\par

``The Book of Blood tells you why the Mountain closed,'' she said. ``The Book of the Beast tells you why it must not open carelessly.''\par

``Because of the Outsiders.''\par

``Because of what happens when the Outsiders find what should remain hidden.''\par

``The Beast.''\par

She looked at me with something almost like pity, controlled and unwelcome. ``Listen to yourself. You say \textit{Beast} as if naming it gives you possession of the fear.''\par

I hated that, because it was true. The word had thrilled me since childhood. Beast. Hell. Chains. Five Keepers. The great last defence, the sacred terror held in reserve for the day the children of the murderer returned to devour the Mountain's shadow. Children love what frightens them, when adults make the fear ceremonial enough. I am telling you now, Paul, across all these years, that I was still that child in the Oasis, and that the word had not yet been made real to me. It would be. There is a room I will tell you about, when you have slept, where the word stopped being a thrill and became a smell.\par

``What does Shield say?'' I asked.\par

``Shield says a thing does not need to be understood by everyone in order to be guarded by everyone.''\par

``That sounds like Ka-Dhavar.''\par

``It should.''\par

``And what do you say?''\par

She was silent long enough that I thought she would not answer. Then one of the wardens near the outer markers shifted, and her eyes moved to him before returning to me. ``I say that a gate is not made holy because there is beauty behind it. Sometimes a gate is holy because there is something terrible behind it, and opening it wrongly would make the opener guilty of every death that followed.''\par

The words were not from a child's lesson. I felt that at once. ``The Beast is real, then.''\par

``I did not say that.''\par

``You almost did.''\par

``No. You almost heard it.''\par

The frustration returned, but now fear entered it. Since Va-Raedin had shown me that a lamp could make a false door more persuasive than a true one, nothing had been clean. The Oasis itself stood before me now as miracle and mechanism, scripture and channel, mercy and ration. Va-Sheva seemed to read enough of that in my face to step back from what she had nearly given me.\par

``The Oasis is smaller than I remembered,'' I said, because retreat had become necessary and I needed ground that did not look like surrender.\par

``Most holy things are.''\par

``And the desert?''\par

``Larger.''\par

I looked toward it. The outer path ran from the Oasis along a fold of stone and vanished between two dark ridges. Beyond it the desert opened in layers of heat and distance. On a rock far out, something flashed briefly, then disappeared. I narrowed my eyes.\par

``What is that?''\par

Va-Sheva's gaze moved instantly. ``Where?''\par

``There. Past the second ridge. The dark stone.''\par

She followed the line. For a moment she said nothing. Then, very softly: ``Do not point.'' My hand had begun to lift. I lowered it. The softness left her face. ``Stay here.'' She turned from me completely.\par

\ornament

\noindent ``A flash,'' he says. ``On a rock out past the second ridge.'' He says it the way a boy reports a small wonder, already moving past it toward the woman and the turned stone, because to him the woman is the story and the flash is only the thing that made her turn her head.\par

I do not move. That is the error, and I know it even as I make it --- not a word, not a wrong word, only a stillness held half a breath too long: the stillness of a man handed the one card he had stopped letting himself pray for. A light, low on the desert, beyond a ridge that does not move, seen from a shade that does not move, at a named hour of the afternoon. It is a line drawn across the sand. I have carried the other line a long time and never the place to cross it. He has just given me the place to stand.\par

He has stopped. ``You have gone still, Paul.''\par

``It is a long account,'' I say, ``and I am old, and the fire is good.'' I let my hand settle on the reels, as though the machine were the only thing in the room that wanted minding. ``The desert throws light like that all over --- mica, old iron, water where no water is. Heat lies to the eye. Your Shield woman knew it. That is why she called it, in the end, probably nothing.''\par

He watches me a breath longer than he has watched me yet. Then he lets it go, because he wants her part of the night more than mine, and because it has not yet crossed him that a servant of record might have a reason of his own to care where a light once stood in a desert. He will remember that I went still. He does not know, tonight, to be afraid of it.\par

``Go on,'' I say.\par

\ornament

\noindent What followed was too quiet to be dramatic, and I have never forgotten it for exactly that reason. There was no shout. No drawn blade. No rush of bodies. Va-Sheva made three gestures in succession: two fingers low, the palm closed, the hand cut toward the outer path. And the Oasis rearranged itself.\par

One warden moved behind the goat pens and became invisible in shade. Another knelt near the dust at the outer marker. A third shifted between the prayer group and the spring. A hidden figure answered from the rock above the descent, where I had seen only stone. Then another signal came from higher on the slope. Shield had been around us the whole time --- visible wardens, hidden bodies, lines of sight, markers arranged so that disturbance could speak before danger had language. I had thought I was standing in a garden. I had been standing in a held breath.\par

The Faith attendant looked up, alarmed, but did not speak. One of the old women tried to rise and was pressed gently back by her companion. A goat stopped chewing. The bees continued.\par

Va-Sheva crouched beside the kneeling warden and touched the ground. She lifted her fingers, rubbed something between them, then looked toward the outer path, the Balance terraces, the prayer group, and the rock where I had seen the flash. A Balance worker had gone still beside the snail trays, water ladle in hand. Va-Sheva pointed to him, then to the shaded wall. He moved at once.\par

The warden near the stones gave a low call. ``Marker.'' The word was ordinary. Shield changed around it. One outer marker had been turned. Barely. A flat side faced inward where the others faced skyward. Va-Sheva did not touch it. She looked at the dust around it, the path beyond, the nearest channel, the position of the prayer group, and the rock beyond the ridge. Then she spoke aloud.\par

``No one returns inside yet.''\par

The Faith attendant's eyes widened. ``The access period---''\par

``Has changed.''\par

``By whose order?''\par

Va-Sheva looked at her. The attendant lowered her head. ``Shield holds the threshold.''\par

``Good.''\par

I heard the words enter the Oasis like a bar lowered across a gate. Faith governed prayer. Balance governed water and work. But here, at the edge of the world, Shield could interrupt sacred timing with two words and no explanation, and everyone bent. That was the first time I understood, in my body rather than my lessons, that the five offices were not equal in every room. In this one, at this line, one of them simply outranked God's own clock.\par

A young Balance boy near the herb beds began to cry. He was old enough to be ashamed of it and too young to stop.\par

Va-Sheva walked the marker line without stepping beyond it. At the turned stone she knelt, took a thin bone pin from her sleeve, and lifted a strand from the ground. It caught the light once, pale against her glove. She placed it in a covered dish handed to her by a warden. ``Second gate,'' she said. A hidden watcher answered from the slope. Then she turned toward the workers.\par

``Count.''\par

An older Balance woman began at once. ``Goat pens: three visible, one shaded, all tied. Snail trays: two open, four covered, no tray missing. Herb knives: five issued, five in hand. Water ladles: eight.''\par

``Show eight.'' Eight ladles lifted. One child's hand shook.\par

``Fibre hooks?''\par

``Three.''\par

``Show three.'' Three hooks lifted.\par

``Outside workers?''\par

``Eleven.''\par

``Name them.'' The Balance woman named them without a slate. One by one, the bodies answered. When all eleven had answered, Va-Sheva nodded to the warden.\par

``Prayer group?'' The Faith attendant named seven. The warden corrected her. ``Six.'' The attendant froze. My stomach tightened before my mind understood.\par

Then one of the old women said sharply, ``Mara.'' The woman beside her looked round. ``She was at the spring rail.'' No one moved.\par

Va-Sheva did. She did not run --- running would have told everyone that fear had outrun command. She crossed the prayer area at a fast walk, two wardens spreading behind her. I saw then that one of the older women from our access group was missing. Not the aunt, not the boy, not the grieving men. An old woman who had knelt nearest the spring basin and moved slowly because one hip did not bend properly. The Faith attendant whispered, ``God keep her in shadow.'' Va-Sheva heard. ``Silence.'' The word cut the air.\par

They found the woman by the spring rail. She had crouched there with one hand against the stone, breathing hard, unable to rise. Whether she had wandered intentionally toward the water or simply lost her balance, I could not tell. In her other hand she held a small wet leaf. A Balance worker made a sound of anger low in her throat. Va-Sheva took the leaf gently, gave it to the worker, then helped the old woman stand. She did not humiliate her. That surprised me. Va-Sheva only leaned close and said something that made the woman close her eyes and nod.\par

Then she signalled. The warden near the channel examined the edge. The turned marker remained under guard. The rock beyond the second ridge flashed no more. The old woman was returned to Faith, who received her with a face drained of colour.\par

Va-Sheva came back to the inner shade. Her eyes passed over the group, counted, and settled briefly on me. ``Seven bodies out,'' she said. ``Seven bodies held. Eleven Balance workers held. No one returns until the marker is reset and the channel is cleared.''\par

The Faith attendant said very quietly, ``She is old.''\par

``Yes.''\par

``She meant no harm.''\par

``Meaning no harm is not the same as remaining harmless.''\par

The old woman began to weep silently, with such shame that I had to look away. The leaf had perhaps been for medicine, or memory, or some private superstition. The reason might have been tender. Shield treated it as a breach until the tenderness could be safely restored. I would learn that grammar very well, later, in rooms where my own tenderness was filed the same way, under the same suspicion, pending the same restoration that did not come.\par

The containment ended as quietly as it had begun. The turned stone was lifted with bone pins and examined beneath a cloth. The channel was checked. A small clay seal was pressed beside the woman's name. The prayer group was recounted. The Balance workers returned to their tasks with the rigid annoyance of people who had lost time but not the permission to complain of it.\par

\ornament

\noindent When Va-Sheva finally came back to me, the official danger had passed, though nothing in her face softened.\par

``You saw the flash first.''\par

``I think so.''\par

``You did not point when told not to.''\par

``No.''\par

``Good.'' The word struck harder than it should have.\par

``Was it connected?'' I asked.\par

``To the marker?''\par

``Yes.''\par

``Probably not.''\par

``Then why---''\par

She looked at me, and I stopped. ``Because two small wrongnesses at a threshold are not treated separately until Shield can afford the luxury,'' she said. ``A flash may be nothing. A turned marker may be nothing. An old woman out of place may be nothing. Three nothings close together become a shape.''\par

\textit{What is unseen can be known by what it disturbs.} \textit{State the consequence of each choice.} \textit{Curiosity is not free.} \textit{Do not admire the flame; ask where the air enters.} All the sentences the Mountain had been pressing into me that week, from Va-Raedin, from Ka-Leth, from the recorder in the repair hall. And now this, from Shield, in the open, over a turned stone: \textit{three nothings close together become a shape.} I understood, dimly, that every office was teaching me the same lesson in its own dialect, and that the lesson was not really about lamps or figs or markers. It was about how to see a thing that did not want to be seen. They were, all of them, training the exact faculty that would one day make me impossible for them to keep.\par

``You see like Light when you say that,'' I said.\par

``No. Light sees what can be shown. Shield sees what must be stopped before it shows itself.'' It was not boastful. That made it more powerful.\par

I looked toward the desert again. The rock beyond the second ridge sat dark and ordinary. The flash might have been mineral, bird, old metal, heat, an eye. ``Have you gone beyond the second ridge?''\par

Va-Sheva did not answer at once. ``That is not a prayer question.''\par

``No.''\par

``Then do not ask it during prayer access.''\par

I nodded. But she had not said no. I have built whole years on smaller silences than that one.\par

\ornament

\noindent The Faith attendant called us back to the prayer stones for the return blessing. The access period had been shortened and lengthened both, made irregular by the containment. Justice would dislike the irregularity. Balance would record the lost task time. Faith would smooth the fear with words. Shield would remember the turned marker longer than anyone remembered the prayer. Each office would keep its own version of that afternoon, and no two versions would agree, and that disagreement was not a failure of the Mountain. It was the Mountain.\par

Va-Sheva stepped aside to let me pass. ``Sa-Kailan.'' I turned. Wind lifted dust between us. With the Mountain shadow behind her and the desert light along her cheek, she seemed less unreachable, because she stood close enough to touch. Then the foolishness of the thought revealed itself. She was not unreachable because of distance.\par

``You did well to notice,'' she said. The praise warmed me before I could defend against it. Then she added, ``But do not start believing the outside speaks to you.'' There it was: the correction inside the gift.\par

``I did not.''\par

``You are fifteen. Everything speaks to you.''\par

``I am not a child.''\par

``No,'' she said. ``Not entirely.''\par

The return blessing required us to wash the dust from our fingers in a shallow bowl before re-entering the descent. I watched the dirt darken the water. The Faith attendant spoke over each of us.\par

``What has touched the open returns under shadow.''\par

We answered. ``Under shadow, we return.''\par

When my turn came I placed my fingers in the bowl. The water was warm now, from sunlight and use. My reflected hand broke around the dust, the image bending, thickening, returning --- the same small miracle that had cost me a night's sleep after the first lesson. But Shield was counting, and Faith was watching, and Balance had already measured the bowl's use, and the outside did not lengthen itself for my curiosity. Even here, even at the edge of the world, I was permitted to see the water bend only for as long as the record allowed.\par

We re-entered the Mountain in single file. The first steps down felt wrong. The passage narrowed quickly. The ceiling returned like a hand over the skull. After the sky, the lamps seemed mean and intelligent. After the desert, the stone smelled of bodies, oil, and old water. Yet relief came too, shameful and immediate. The Mountain enclosed me, and some part of me was grateful to be enclosed. Outside had not been freedom. Remember that, Paul, when your century tells you the open air is always liberty. Sometimes the open air is only a larger place to be watched from.\par

At the second threshold I returned the clay counter. The Balance clerk checked the number against the slate. Shield marked seven bodies in. Faith took the dust bowl away. A Justice witness, seated almost invisibly near the wall, made a mark after the gate closed. The cedar face of the inner door shut. The iron bar fell. The Mountain informed itself that we had returned.\par

I stood in the passage a moment too long. Behind me lay the Oasis: water in shadow, bees, goats, snails, herbs, channels, prayer stones, the Book of Blood spoken under open sky, and Va-Sheva's hand cutting the air before an explanation could catch up to it. Beyond it lay the desert, immense and bright and watched from places I had not known to look. Before me lay the corridors of the Mountain, lit by measure, closed by rule, alive with permissions.\par

The Oasis was proof. I could not deny that. A hidden spring in the Mountain's shadow, water held against the open world, food and fibre and medicine and wax and milk and seal and wick and hide and cord and prayer and count. It made the Covenant physical in a way no gravure could. But the reading had made the desert proof too. It was where the younger brother had fled. It was where holy speech had been stolen and broken. It was where the children of the murderer had wandered under the eye of Heaven, remembering the Mountain without finding it, twisting the Gospel into hunger. The desert was not freedom in the Book of Blood. It was inheritance. It was curse.\par

And the Book of the Beast had added something worse, though Va-Sheva had barely opened the door to it. If Blood explained why the Mountain had closed, Beast explained why the closing could not fail. Five Keepers. Five charges. Five hands kept apart because whatever slept beneath the Mountain was not merely a story for frightened children. Or, if it was a story, it was a story Shield treated like a gate.\par

I touched the folded permission bark inside my robe, though it had already spent its authority. I want to end the day here, Paul, on that small useless gesture, because it is the truest thing I did that afternoon. The outside had not released me from the Mountain. It had only explained why the Mountain held so tightly. A boy went out to see the sky and came back understanding his cage better, and mistook the understanding for a kind of freedom. I would make that mistake many more times before I learned that a clearer view of the bars is not the same as an open door. But I was fifteen, and the water had bent for me once, out there, in the light no office owned, and I fell asleep that night certain the desert had flashed for me alone.\par

\ornament

\noindent ``Flashed for me alone,'' he says, and very nearly smiles at the boy he is describing, the one who believed a whole desert could pick him out of a group of seven.\par

I let the smile stand. I do not tell him that the boy was not entirely wrong --- that the light did have a witness it was meant to reach, and that the witness sits across the table from him now, a long lifetime late, holding a reel that turns. He thinks he has given me a boy's evening: a woman he wanted and could not have, a sky too large for him, a stone turned by nothing. He has given me a line across a desert, and a place to stand to read it, and the hour the shadow falls true.\par

And he believes he is telling it home. That is the part I keep for myself and do not set on the table between us. He gives it freely, thanks me for the listening, and means every word of it for his own. I let him. The reels turn. Tomorrow he will hand me a little more of the only thing my house still wants, and call it his account, and sleep well after.\par

``Sleep,'' I say. ``You have given enough for one night.''\par

He does not know how much that is. Nor, quite, do I --- only that I am nearer it tonight than I was this morning, by one light on a far rock; and the boy who saw it first is the man who has just handed it away.\par

\ornament
